\begin{frame}{Выводы}
	\small
	\begin{columns}[T] % Выравнивание колонок по верхнему краю
		
		\begin{column}{0.48\textwidth}
			\begin{itemize}
				\item Применён численный подход к разрешению систем условий порядка с помощью \texttt{scipy.optimize} вместо символьных вычислений.
				
				\item Алгоритм роя частиц дополнен локальной полировкой Нелдера--Мида.
				
				\item Достигнуто снижение глобальной погрешности:
				\begin{itemize}
					\item осциллятор --- в \textbf{6842} раза;
					\item уравнение Шрёдингера --- в \textbf{92.52} раза;
					\item задача двух тел --- в \textbf{7.3} раза.
				\end{itemize}
			\end{itemize}
		\end{column}
		

		\begin{column}{0.48\textwidth}
			\begin{itemize}
				\item Минимизация фазовой и диссипативной ошибок позволила обеспечить долговременную устойчивость физических инвариантов при интегрировании.

				\item Интеграция нейросетевых аппроксиматоров для замещения численного решателя ОДУ позволит кратно снизить вычислительные затраты при расчёте целевой функции для задач высокой размерности.
			\end{itemize}
		\end{column}
		
	\end{columns}
\end{frame}